\chapter{Otitis media}
\index{Otitis media}

\medskip
\noindent\textit{Educational information; seek professional advice for individual care.}

\section*{What it is}
Otitis media is inflammation or infection in the middle ear, the space behind the eardrum. It commonly follows a cold and is especially common in young children. Fluid can remain after the infection and temporarily reduce hearing.

\section*{Symptoms}
Symptoms may include ear pain, fever, reduced hearing, pressure or fullness, irritability, poor sleep, or fluid draining from the ear if the eardrum has perforated. Babies may pull at the ear, feed poorly, or seem unsettled. Symptoms can also be caused by outer-ear infection, wax, dental problems, or other conditions.

\section*{Assessment and treatment}
Diagnosis usually requires looking at the eardrum and considering the person’s age, severity, duration, and general health. Many uncomplicated infections improve without antibiotics. Pain relief and fluids may be enough; antibiotics or specialist review may be needed for selected children, severe illness, persistent symptoms, repeated infections, or complications. Do not put drops or objects in the ear unless advised.

\section*{Hearing and follow-up}
Arrange follow-up if hearing remains reduced, fluid persists, speech or learning is affected, or infections recur. A clinician may monitor for glue ear and consider hearing assessment or ear-nose-throat referral.

\section*{When to seek urgent care}
Seek urgent help for swelling or redness behind the ear, a protruding ear, severe headache, neck stiffness, confusion, facial weakness, severe dizziness, or a very unwell child. Arrange medical review for persistent fever, worsening pain, new discharge, or symptoms in a very young infant.

\section*{Sources}
\begin{itemize}
\item NHS ear infection guidance: \url{https://www.nhs.uk/conditions/ear-infections/}
\end{itemize}
